\input{../preamble.tex}
\SetDate[13/05/2026]
\setbeamertemplate{page number in head/foot}{}

\begin{document}

\section{Automatismes}

\begin{frame}

\centering \huge
Automatismes

\large
Calculatrice autorisée

\end{frame}

\subsection{Moyennes}

\framedelayed{
	Calculer la moyenne arrondie à $10^{-2}$ près de la série statistique suivante.
	\boxAB{
		\[ X= (0 ; -1 ; 12 ; -11 ; -11 ; 3 ; 0). \]
	}{
		\[ X= (3 ; 2 ; 15 ; -8 ; -8 ; 6 ; 3). \]
	}
}{
	\boxAB{
		$X_A = (0 ; -1 ; 12 ; -11 ; -11 ; 3 ; 0)$
		$\overline{X_A} = \frac{0+(-1)+12+(-11)+(-11)+3+0}{7} \approx -1,14$
	}{
		
		$X_B = (3 ; 2 ; 15 ; -8 ; -8 ; 6 ; 3) $
		$\overline{X_B} = \frac{3+(2)+15+(-8)+(-8)+6+3}{7} \approx 1,86$
	}
	Remarque : $X_B = X_A + 3$ donc $\overline{X_B} = \overline{X_A} + 3$ par linéarité de la moyenne.
}


\framedelayed{
	Calculer la moyenne arrondie à $10^{-2}$ près de la série statistique suivante.
	
\def\arraystretch{1.5}
\setlength\tabcolsep{15pt}
	\boxAB{
		\begin{tabular}{|c|c|c|c|c|}\hline
		Valeur & $-2,5$ & $-4,5$ & $15,1$ & $9,825$  \\ \hline
		Effectif & 3 & 2 & 10 & 50 \\ \hline
		\end{tabular}
	}{
		\begin{tabular}{|c|c|c|c|c|}\hline
		Valeur & 5 & 9 & $-30,2$ & $-19,65$ \\ \hline
		Effectif & 3 & 2 & 10 & 50 \\ \hline
		\end{tabular}
	}
}{
\def\arraystretch{.5}
\setlength\tabcolsep{10pt}
	\boxAB{
		\begin{tabular}{|c|c|c|c|c|}\hline
		Valeur & -2,5 & -4,5 & 15,1 & 9,825  \\ \hline
		Effectif & 3 & 2 & 10 & 50 \\ \hline
		\end{tabular}
		$\overline{X_A} = \frac{-2,5\times3 -4,5\times2 +15,1\times10 + 9,825\times50}{3+2+10+50} \approx 9,62. $
	}{
		\begin{tabular}{|c|c|c|c|c|}\hline
		Valeur & 5 & 9 & $-30,2$ & $-19,65$ \\ \hline
		Effectif & 3 & 2 & 10 & 50 \\ \hline
		\end{tabular}
		$\overline{X_B} = \frac{5\times3 + 9\times2 - 30,2\times10 - 19,65\times50}{3+2+10+50} \approx -19,25.$
	}
	Remarque : $X_B = -2 \times X_A$ donc $\overline{X_B} = -2\times\overline{X_A}$ par linéarité de la moyenne.
}


\subsection{Écart-type}

\framedelayed{
	Soient $\sigma_1$ l'écart-type de la série $X_1$ et $\sigma_2$ l'écart-type de la série $X_2$.
	
	A-t-on $\sigma_1 < \sigma_2$ ou $\sigma_1 > \sigma_2$ ?
	\boxAB{
		$X_1 = (10 ; 9 ; 11)$ et $X_2 = (6 ; 11 ; 13)$.
	}{
		$X_1 = (4 ; 12 ; 14)$ et $X_2 = (10 ; 9 ; 11)$.
	}
}{
	\boxAB{
		$X_1 = (10 ; 9 ; 11)$ et $X_2 = (6 ; 11 ; 13)$.
		
		Les moyennes sont toutes les deux 10, donc
		$\sigma_1 < \sigma_2$ car la première série est plus concentrée autour \mbox{de 10}.
	}{
		$X_1 = (4 ; 12 ; 14)$ et $X_2 = (10 ; 9 ; 11)$.
		
		Les moyennes sont toutes les deux 10, donc
		$\sigma_1 > \sigma_2$ car la première série est plus dispersée autour de 10.
	}
}

\subsection{Quartile}

\framedelayed{
	Donner l'effectif total de la série statistique décrite par le diagramme à barres ci-dessous.

	\begin{center}
	\begin{tikzpicture}[scale=1]
	\begin{axis}[
		xmin=0, xmax=15,
		ymin=0, ymax=8,
		x=18pt,
		y=10pt,
		grid=both,
		grid style = {opacity=.5},
		xtick distance=1,
		axis lines=center,
		axis line style = -,
		extra x ticks={0},
		extra y ticks={0},
		xlabel=Valeur,
		xlabel near ticks,
		ylabel=Effectif,
		ylabel near ticks,
		ybar,
		bar width=8pt,
	]     
      \addplot+[bar shift=0,draw=BLACK, thick, fill=RED_E!50] plot coordinates {
        (3,8)
        (4,6)
        (5,2)
        (9,4)
        (10,3)
        (12,5)
        (13,2)
      };
      
      \addplot+[bar shift=0,draw=BLACK, thick, fill=GREEN_E!50] plot coordinates {
        (6,8)
        (8,3)
        (11,3)
        (14,5)
        (2,6)
        (1,2)
      };
	\end{axis}
	\end{tikzpicture}
	{\color{RED_E} En rouge : série A}
	\hspace{2cm}
	{\color{GREEN_E} En vert : série B}
	\end{center}
	
}{
	\begin{center}
	\begin{tikzpicture}[scale=1]
	\begin{axis}[
		xmin=0, xmax=15,
		ymin=0, ymax=8,
		x=18pt,
		y=9pt,
		grid=both,
		grid style = {opacity=.5},
		xtick distance=1,
		axis lines=center,
		axis line style = -,
		extra x ticks={0},
		extra y ticks={0},
		xlabel=Valeur,
		xlabel near ticks,
		ylabel=Effectif,
		ylabel near ticks,
		ybar,
		bar width=8pt,
	]     
      \addplot+[bar shift=0,draw=BLACK, thick, fill=RED_E!50] plot coordinates {
        (3,8)
        (4,6)
        (5,2)
        (9,4)
        (10,3)
        (12,5)
        (13,2)
      };
      
      \addplot+[bar shift=0,draw=BLACK, thick, fill=GREEN_E!50] plot coordinates {
        (6,8)
        (8,3)
        (11,3)
        (14,5)
        (2,6)
        (1,2)
      };
	\end{axis}
	\end{tikzpicture}
	\end{center}
	\boxAB{
		$8+6+2+4+3+5+2 = 30$
	}{
		$2+6+8+3+3+5 = 27$
	}
}



\section{Solutions}

\shipoutAnswer

\end{document}